%% notes-tex/024-perturb-seq-pilot/sections/4-open.tex
%% SOURCE: hand-authored synthesis. The environment arithmetic is
%%         experiments/024-perturb-seq-costing/scripts/scaling_analysis.py
%%         (Sec. 7.3 of the costing review) and pilot_options.py (Fig. 1b); the
%%         delivery facts are Sec. 7.1 of the same review. No new numbers.

\section{The Two Decisions That Were Open\secstatus{todo}}
\label{sec:open}

\subsection{Growth Media, or the Readout First\secstatus{todo}}
\label{sec:media}

Both, and the split falls along the two hosts rather than along a preference.

Growth media are the cheap axis. Cell demand grows in proportion to the number
of conditions, where covering named pairs of genes grows with the square of the
panel, and any route whose first barcode is a plate of wells labels conditions
for nothing (\cref{fig:pilots}b). A model that proposes strain designs also
needs both halves of its input to vary, and a dataset that varies only genotype
in one medium constrains only the genotype half.

What media cannot do is multiply an assay that does not exist yet. Running four
conditions through a readout whose guide detection is unmeasured produces four
copies of the same unanswered question, at four times the price. So the ordering
is not a judgment about which is more interesting.

\begin{itemize}[leftmargin=1.4em,itemsep=2pt,topsep=2pt]
  \item \textbf{\org{S.\ cerevisiae} spends its first run on the readout}, in
    one condition, because the library exists and $q$ is the gate
    (\cref{sec:guide}). Media come second, and by then the plate route is the
    way to buy them.
  \item \textbf{\org{P.\ kudriavzevii} spends its first run on media}, because it
    has no library to gate on. A single unperturbed condition under-uses a
    channel that has already been paid for, and the honest way to thicken that
    run is more conditions rather than genetics that do not exist yet.
\end{itemize}

On the plate route eight media ride in one protocol run for \$4{,}990 plus the
one-time plate purchase, against \$12{,}550 for four media as four droplet
channels (\cref{tab:pilots}). That comparison is the argument for the plate, and
it only applies once the plate chemistry is known to work on this cell wall,
which nobody has shown. Running the \org{P.\ kudriavzevii} first experiment on
droplet across a few media, rather than on the plate across many, buys the wall
answer on a platform the core supports and leaves the plate for run two.

\subsection{Plasmid Copy Number\secstatus{todo}}
\label{sec:copy}

Not early, with one exception worth building into the first run.

\notesec{Copy number is not guide diversity, and conflating them is the trap} A
cell that takes up one plasmid molecule and amplifies it to fifty copies carries
fifty copies of \emph{one} guide. Diversity comes only from the number of
independent uptake events at transformation. Raising average copy number by
weakening the selection marker is well established and finely tunable in yeast,
over ``5--100 copies per
cell''~\citep{lianConstructionPlasmidsTunable2016}, and it does not by itself
put two guides in a cell.

\notesec{Three reasons to keep the pilot low-copy} A single-guide screen is
interpretable precisely because one cell means one guide, and the field builds
its libraries in ``typically, low-copy number vectors to prevent multiple guide
plasmids in a single
cell''~\citep{robertsonAdvancesCRISPRenabledGenomewide2025}. Multiple plasmids
per cell are also transient: multi-vector transformants segregate over roughly a
day of outgrowth~\citep{scanlonQuantifyingResolvingMultiple2009}, and 2$\mu$
plasmids are retained in only about 60\% of cells even in a laboratory
strain~\citep{haysNaturalVariantEssential2020}, so a high-copy pilot measures a
population that is changing while it is being measured. And a cell carrying two
guides at unknown copy number confounds the one number the first run exists to
produce, since a guide missed in such a cell and a guide missed in a
single-guide cell are not the same failure.

\notesec{The exception, and it is a hypothesis} Guide detection depends on how
much guide transcript is present, and copy number raises that.
\textbf{Hypothesis (untested):} a higher-copy vector raises $q$ by putting more
guide RNA in each cell. Nothing measures this, and it is directly testable in
the run already planned, as a second channel carrying the same library on a
higher-copy backbone (\cref{tab:pilots}, row two, \$8{,}576 for the pair rather
than \$4{,}493). Worth doing if the budget allows, and worth skipping if the
first channel returns a high $q$ on its own.

\notesec{When copy number does become the right lever} Once $q$ is known and
single-guide screening works, putting several guides in one cell is what makes
combinations affordable, and the multi-plasmid route is the only way past the
two-to-four guide ceiling that cloning arrays into one vector imposes. The
missing measurement then is the per-cell distribution of distinct guides under
tuned delivery, which the literature does not report for yeast. That is the
second pilot, not the first.

\begin{keybox}
\textbf{Keep the first library low-copy, and treat high copy as a detection
experiment rather than a multiplexing one.} Copy number buys guide transcript
per cell, not guides per cell, so it cannot deliver combinations and it can
confound the measurement of $q$. The one version worth running now is a second
channel of the same library on a higher-copy backbone, purely to see whether $q$
rises with it.
\end{keybox}
